\documentclass[10.5pt]{article}
\usepackage[letterpaper, margin=0.6in, top=0.5in, bottom=0.5in]{geometry}
\usepackage{amsmath,amssymb,amsthm}
\usepackage{fancyhdr}
\usepackage{xcolor}
\usepackage{tikz}
\usepackage{enumitem}
\usepackage{tcolorbox}
\usepackage{hyperref}

% Define colors matching instructor board annotations
\definecolor{chalkblue}{RGB}{0, 105, 185}
\definecolor{chalkred}{RGB}{205, 35, 35}
\definecolor{chalkpink}{RGB}{230, 90, 120}
\definecolor{chalkcyan}{RGB}{15, 160, 190}
\definecolor{chalkgreen}{RGB}{20, 130, 50}
\definecolor{boxbg}{RGB}{252, 252, 254}
\definecolor{boxframe}{RGB}{150, 160, 180}

\pagestyle{fancy}
\fancyhf{}
\renewcommand{\headrulewidth}{0pt}
\renewcommand{\footrulewidth}{0pt}

\setlength{\parindent}{0pt}
\setlength{\parskip}{2.5pt}

\begin{document}

% ==============================================================================
% PAGE 1
% ==============================================================================
\noindent \textbf{\large Math 105L} \hfill \textbf{\large Lab 5 Part 2}\\[1pt]
{\textbf{\Large Fundamental Differentiation Rules \& Negative Exponents}}\\[-5pt]
\rule{\textwidth}{1.2pt}

\vspace{0.03in}

\textbf{\large The Three Core Differentiation Rules}

In this laboratory demonstration, we synthesize the foundational differentiation building blocks:

\begin{tcolorbox}[colback=boxbg, colframe=boxframe, boxrule=0.8pt, arc=3pt, left=8pt, right=8pt, top=4pt, bottom=4pt]
\begin{itemize}[leftmargin=1.2em, topsep=1pt, itemsep=3pt]
  \item \textbf{Sum Rule:} The derivative of a sum is the sum of the derivatives:
  \[
    \mathbf{\frac{d}{dx}\big(f(x) + g(x)\big) = f'(x) + g'(x)}
  \]
  \item \textbf{Constant Multiple Rule:} Constants factor cleanly out of differentiation:
  \[
    \mathbf{\frac{d}{dx}\big[C \cdot f(x)\big] = C \cdot f'(x)}
  \]
  \item \textbf{Power Rule:} For any real exponent $n$:
  \[
    \mathbf{\frac{d}{dx}\big[x^n\big] = n x^{n-1}}
  \]
\end{itemize}
\end{tcolorbox}

\vspace{0.04in}

\textbf{\large Comprehensive Example Worked on Blackboard: $f(x) = \frac{6}{x} + 7x^{10}$}

We differentiate $f(x)$ with full pedagogical justification at every algebraic step:

\begin{itemize}[leftmargin=1.2em, topsep=1pt, itemsep=4pt]
  \item \textbf{Step 1: Rewrite reciprocals as negative powers:}
  \[
    f(x) = \frac{6}{x} + 7x^{10} = \mathbf{6x^{-1} + 7x^{10}}
  \]
  \item \textbf{Step 2: Apply the Sum Rule:}
  \[
    f'(x) = \frac{d}{dx}\big[6x^{-1}\big] + \frac{d}{dx}\big[7x^{10}\big] \quad \text{\footnotesize (Sum Rule)}
  \]
  \item \textbf{Step 3: Apply the Constant Multiple Rule (twice):}
  \[
    f'(x) = 6 \cdot \frac{d}{dx}\big[x^{-1}\big] + 7 \cdot \frac{d}{dx}\big[x^{10}\big] \quad \text{\footnotesize (Constant Multiple Rule)}
  \]
  \item \textbf{Step 4: Apply the Power Rule (twice):}
  \begin{itemize}[label=--]
    \item For $x^{-1}$: $n = -1 \implies \frac{d}{dx}[x^{-1}] = (-1) x^{-1 - 1} = -x^{-2}$.
    \item For $x^{10}$: $n = 10 \implies \frac{d}{dx}[x^{10}] = 10 x^{10 - 1} = 10 x^9$.
  \end{itemize}
  \[
    f'(x) = 6\big(-x^{-2}\big) + 7\big(10 x^9\big) \quad \text{\footnotesize (Power Rule)}
  \]
  \item \textbf{Step 5: Simplify and express with positive exponents:}
  \[
    \mathbf{f'(x) = -\frac{6}{x^2} + 70x^9} \qquad \checkmark
  \]
\end{itemize}

\vspace{0.03in}

\begin{tcolorbox}[colback=chalkblue!6, colframe=chalkblue, boxrule=1pt, arc=3pt, left=8pt, right=8pt, top=4pt, bottom=4pt]
\textbf{Key Takeaway for Fractions and Polynomials:}\\[1pt]
Before applying the Quotient Rule to expressions with constant numerators like $\frac{C}{x^k}$, always rewrite them algebraically as $C x^{-k}$. Applying the Power Rule to negative exponents is much faster and significantly less error-prone!
\end{tcolorbox}

\end{document}

